\documentclass[../../main.tex]{subfiles}
\begin{document}

\begin{center}
  {\fontsize{48pt}{58pt}\selectfont \textbf{1999}\normalsize\textbf{年度}}
\end{center}

\begin{table}[htb]
  \centering
  \caption{出題テーマと難易度}
  \begin{tabular}{|c|p{9cm}|c|}
    \hline
    \textbf{問題番号} & \textbf{テーマ} & \textbf{難易度} \\
    \hline
    第1問 & 放物線の弦の中点の軌跡 & \\
    \hline
    第2問 & ベクトルの内積条件から定まる楕円 & \\
    \hline
    第3問 & 並べ替え不等式と和の下界 & \\
    \hline
    第4問 & 複素数で表す正三角形 & \\
    \hline
    第5問 & 平方根の有理独立性 & \\
    \hline
    第6問 & 媒介変数表示の曲線と面積 & \\
    \hline
  \end{tabular}
\end{table}

\begin{table}[htb]
  \centering
  \caption{解答}
  \renewcommand{\arraystretch}{1.4}
  \begin{tabular}{|c|c|p{9cm}|}
    \hline
    \textbf{問題番号} & \textbf{小問} & \textbf{答え} \\
    \hline
    第1問 & -- & $y=x^2+\dfrac{6^{2/3}}{4}$ \\
    \hline
    第2問 & -- & $AB=2t$として，$A,B$の中点を原点，$AB$を$x$軸にとると$\dfrac{x^2}{t^2+c/2}+\dfrac{y^2}{c/2}=1$ \\
    \hline
    第3問 & (1) & 指定された不等式が成立（証明問題） \\
    \cline{2-3}
    & (2) & $\sum_{k=1}^n\dfrac{x_k^2}{k^2+1}>n-\dfrac{8}{5}$（証明問題） \\
    \hline
    第4問 & (1) & $\{\beta,\gamma\}=\{1+\omega z,1+\omega^2z\}$，$\omega=e^{2\pi i/3}$ \\
    \cline{2-3}
    & (2) & $(\arg\alpha,\arg\beta,\arg\gamma)=(\pi/9,4\pi/9,16\pi/9)$ \\
    \hline
    第5問 & (1) & $p+q\sqrt{2}+r\sqrt{3}=0\implies p=q=r=0$（証明問題） \\
    \cline{2-3}
    & (2) & $f(1),f(1+\sqrt{2}),f(\sqrt{3})$のいずれかは無理数（証明問題） \\
    \hline
    第6問 & -- & $0\le x\le1$，$0\le y\le3(2\sqrt{3}-3)$．面積は$28-40\log2$ \\
    \hline
  \end{tabular}
\end{table}

\paragraph{試験形態}
6問．試験時間は要確認．

\end{document}
